% !TeX root=../main.tex
% در این فایل، عنوان پایان‌نامه، مشخصات خود و چکیده پایان‌نامه را به انگلیسی، وارد کنید.

%%%%%%%%%%%%%%%%%%%%%%%%%%%%%%%%%%%%
\latinuniversity{University of Tehran}
\latincollege{College of Engineering}
\latinfaculty{Faculty of Engineering Science}
\latindepartment{Control}
\latinsubject{Electrical Engineering}
\latinfield{Control}
\latintitle{Design and Fabrication of a Soft Tactile Sensor based on sensory data fusion with application in Object Manipulation}
\firstlatinsupervisor{Dr. Mehdi Tale Masouleh}
\secondlatinsupervisor{Dr. Ahmad Kalhor}
\firstlatinadvisor{Dr. Mohammad Reza Nayeri}
%\secondlatinadvisor{Second Advisor}
\latinname{Mohammad Reza}
\latinsurname{SheykhAzimi PoorSardroud}
\latinthesisdate{Septemmber 2025}
\latinkeywords{Multi-modal, Bio-inspired, Tactile sensor, Grasp force}
\en-abstract{
This research presents the design, fabrication, calibration, and evaluation of a soft, multi-modal, bio-inspired tactile sensor with the aim of enriching the interactive capabilities of robots. Drawing inspiration from the multi-modal nature of human touch, which perceives static and dynamic information with different receptors, this thesis addresses a novel approach to overcome the limitations of single-modality sensors. The research path began with the design and fabrication of a single-axis prototype based on a barometric pressure sensor; while proving the mechanism's efficacy in accurately measuring normal force, its linearity and hysteresis properties were also assessed. As this prototype lacked the capability for multi-modal sensing and tangential force measurement, an integrated multi-modal sensor was developed. This sensor utilizes a barometric sensor, two Hall effect sensors with an embedded magnet to resolve tangential forces, and a MEMS microphone—acting as a piezoelectric sensor—to capture high-frequency vibrations from slip and surface texture. To manage and synchronize the large volume of heterogeneous data, a custom data acquisition board was designed with innovative features like a composite USB communication protocol and advanced utilization of Direct Memory Access (DMA). Subsequently, due to the sensor's inherently nonlinear and complex nature, the calibration process was performed using deep learning models. A fully connected neural network was employed to model the static behavior and accurately extract force vectors, while a Long Short-Term Memory (LSTM) network was used to characterize the sensor's dynamic and time-dependent properties, with evaluation results confirming the high accuracy. Finally, the overall system's efficacy was demonstrated in a practical implementation. The sensor was mounted on a robotic gripper, and within a control loop, its ability for precise grasp force control while concurrently detecting dynamic phenomena based on vibrational data was successfully showcased. These achievements provide a comprehensive hardware and software framework for advancing robotic tactile perception and achieving intelligent and safe interactions with objects.
}
